\section{Threats to Validity}
\label{sec:threats}

\SecNote{States the limits plainly so that reviewers do not have to find them. Carried by the theory of Section II-B and by reading the code. Several threats can be removed by re-running; those carry a \Fix{} and may shrink.}

\textbf{Exchangeability and what the bound means.} Conformal risk control assumes exchangeable data, and CI builds are not exchangeable, since they are ordered and code, tests and failure rates change over time. Exponential weights favor recent builds and follow the setting of Barber et al.~\cite{barber2023beyond} and Farinhas et al.~\cite{farinhas2024nonexchangeable}, whose results bound the degradation by a term in total variation distances that we cannot estimate. Any bound we could claim is marginal. It holds on average across builds, not for each build, and a single build may miss all its failures while the average stays within target. The bound does not apply to the combined selector or to the LLM judge, because neither was part of the calibration. For these reasons we report measured recall and do not present calibration as an assurance of it \Fix{N1, N15}.

\textbf{Calibrated versus reported quantity.} The loss averages over builds, with zero loss on passing builds, whereas fault recall pools over failures. These differ whenever failures are rare or concentrated, and the difference is visible in RQ1 \Fix{N16, N17}.

\textbf{Non-stationarity.} ShardingSphere's slices have 900, 208 and 239 failures. Thresholds selected on one regime may not carry over to the next, and the chronological split gives one such transition per project \Fix{N16: a rolling-origin evaluation would give several}.

\textbf{Temporal leakage in the index.} Test summaries are generated from the latest version of each test file in the studied history, which postdates many calibration and evaluation builds. Scores may therefore benefit from information a deployment would not have, which could inflate recall and affects the case studies too \Fix{N14: regenerate point-in-time summaries, or measure the effect on a sample}.

\textbf{Failure labels.} A failure in this benchmark is a failed test file in a build. It is not a confirmed regression, and we have no flakiness labels. The guard's results suggest that repeated failures dominate ShardingSphere's stream, and some of these may be flaky. Predictive selection at Facebook found flaky outcomes far outnumbered failures after retries and warned that models trained on them mislead~\cite{machalica2019predictive}. Fault recall in this paper should be read as recall of observed test failures \Fix{N30: rerun or classify the 18 at minimum}. For the same reason, failure counts at file level may overstate independent faults; one faulty commit can produce several failures, and the 239 failures come from 34 builds \Fix{N31: change recall, intervals}.

\textbf{Feature provenance.} The guard and the ranker use precomputed benchmark features. We rely on these being computed only from builds before the current one, as the benchmark's definitions imply~\cite{yaraghi2023tcpci} \Fix{N38: verify on data}. The High-Savings thresholds are fixed in the code, with no recorded selection procedure \Fix{N15}.

\textbf{Cost constructs.} Net savings rest on an assumed 5\,s per test, serial execution, and list prices for runners. They omit the cost of a missed failure, index refresh, and feature collection for the ranker. Carbon estimates use fixed constants and exclude inference energy \Fix{N19, N22, N23, N24}. All dollar values come from one run over a window of several months and should be read as orders of magnitude.

\textbf{Implementation and reproducibility.} LLM outputs are nondeterministic and depend on the model version, and we report one cached run \Fix{N31, N37}. Diffs, sources and prompts are truncated, which may hide relevant changes in large commits \Fix{N36}. The reported runs need to be regenerated from a single commit of the pipeline \Fix{N34}.

\textbf{External validity.} We studied two open-source Java projects from one benchmark. Log4j2, with its two failures, is a case study with few failures and supports no claim about generalization, and its threshold was not calibrated on its own failures \Fix{N3}. Effectively, quantitative findings rest on one project and 34 failing builds \Fix{N35}. A company's monorepo, its failure rate and its test durations may behave differently.
